  \def\stat{zac-suchkov}



\def\tit{УГРОЗЫ И РИСКИ РЕАЛИЗАЦИИ КОМПЛЕКСНЫХ НАУЧНО-ТЕХНИЧЕСКИХ 

ПРОГРАММ\\ В~РАМКАХ ПРИОРИТЕТОВ СТРАТЕГИИ\\ НАУЧНО-ТЕХНОЛОГИЧЕСКОГО 

РАЗВИТИЯ РОССИИ$^*$}



\def\titkol{Угрозы и~риски реализации КНТН %комплексных научно-технических программ 

в~рамках приоритетов Стратегии} % научно-технологического  развития России}





\def\aut{А.\,А.~Зацаринный$^1$, А.\,П.~Сучков$^2$}



\def\autkol{А.\,А.~Зацаринный, А.\,П.~Сучков}



\titel{\tit}{\aut}{\autkol}{\titkol}



\index{Зацаринный А.\,А.}

\index{Сучков А.\,П.}

\index{Zatsarinny A.\,A.}

\index{Suchkov A.\,P.}







{\renewcommand{\thefootnote}{\fnsymbol{footnote}}

\footnotetext[1]

{Статья подготовлена при поддержке РФФИ (проект 18-29-03091).}} 





\renewcommand{\thefootnote}{\arabic{footnote}}

\footnotetext[1]{Институт проблем информатики Федерального исследовательского центра <<Информатика и~управление>> 

Российской академии наук, \mbox{AZatsarinny@ipiran.ru}}

\footnotetext[2]{Институт проблем информатики Федерального исследовательского центра <<Информатика и~управление>> Российской академии наук, \mbox{ASuchkov@ipiran.ru}}



 \label{st\stat}



 \thispagestyle{headings}

 

 \vspace*{-12pt}  



  

  



  \Abst{Рассмотрены актуальные проблемы реализации Стратегии научно-технологического 

  развития России в~рамках принятых приоритетов на основе 

комплексных на\-уч\-но-тех\-ни\-че\-ских программ исследований и~разработок (КНТП), 

прежде всего в~рамках первого приоритета. Приведены методические подходы 

к~классификации рисков и~угроз на основе их систематизации, которые препятствуют 

выполнению принятых КНТП. Предложен классификатор угроз, позволяющий получить 

некоторое цельное представление об эффективности запланированного в~КНТП 

комплекса мероприятий. Дана оценка потенциальных рисков и~угроз применительно 

к~КНТП <<Искусственный интеллект как драйвер цифровой трансформации экономики 

России>>, разработанной Федеральным исследовательским центром <<Информатика 

и~управ\-ле\-ние>> РАН. Показано, что основные риски обусловлены недостаточным 

стимулированием научных исследований и~слабой востребованностью инновационных 

научных результатов.}



  \KW{комплексная научно-техническая программа; стратегия на\-уч\-но-тех\-но\-ло\-ги\-че\-ско\-го 

развития; классификация рисков и~угроз; искусственный интеллект; научные 

исследования}



\DOI{10.14357\!/\!08696527200309}



\vskip 10pt plus 9pt minus 6pt



 \thispagestyle{headings}

 

 %\vspace*{-12pt}



\section{Введение}



    В целях обеспечения реализации приоритетов, определенных Стратегией 

на\-уч\-но-тех\-но\-ло\-ги\-че\-ско\-го развития Российской Федерации~[1], 

Правительство РФ приняло Постановление, в~котором утверждаются Правила 

разработки, утверждения, реализации, корректировки и~завершения 

КНТП полного 

инновационного цикла~[2]. Комплексная на\-уч\-но-тех\-ни\-че\-ская программа определена как совокупность 

скоординированных по задачам, срокам и~ресурсам мероприятий, 

включающих в~себя научные исследования и~этапы инновационного цикла до 

создания технологий, продукции и~оказания услуг. Предполагается, что 

КНТП формируются для создания прорывных отечественных технологий 

и~получения результатов, обеспечивающих повышение 

конкурентоспособности экономики. Другими словами, КНТП 

рассматриваются как один из действенных инструментов для выполнения 

амбициозной задачи вхождения России в~пятерку ведущих экономик мира.

    

    Заметим, что эта задача представляется весьма сложной, учитывая 

нынешнее состояние экономики России, которая в~рейтинге 2019~г.\ 

занимает только 11-е место~[3]. Реально войти в~пятерку ведущих стран (ее 

в~2019~г.\ составили США, Китай, Япония, Германия и~Великобритания) 

Россия может только на основе повышения эффективности научных 

исследований, быстрой инновационной реализации результатов и,~как 

следствие, увеличения доли высокотехнологичной продукции 

в~производстве. При этом в~Стратегии отмечено, что в~России крайне низкие 

показатели в~этой части (доля инновационной продукции всего 8\%--9\%, 

доля экспорта высокотехнологичной продукции в~мировом объеме всего 

0,4\%)~[1]. 

{\looseness=1



}

    

    В связи с~этим введение на государственном уровне такого инструмента, 

как КНТП, представляется важным и~своевременным. 

    

    Вместе с~тем, чтобы обозначенный инструмент в~форме КНТП был 

приведен в~действие, необходима эффективная организация работ. Так, 

в~соответствии с~Постановлением должна быть определена организационная 

система субъектов процесса реализации КНТП (инициатор, ответственный 

исполнитель, соисполнитель, заказчик, участники, базовая организация), 

а~также весьма непростая организационная схема их взаимодействия. 

    

    В связи с~этим представляется актуальным рассмотрение и~анализ 

потенциальных угроз и~рисков в~рамках осуществления указанных 

процессов. В~статье предложены методические подходы к~созданию 

классификации таких угроз на основе их систематизации с~целью получения 

некоторого цельного представления об эффективности запланированного 

в~КНТП комплекса работ согласно замыслу Постановления~[2]. 

    

    Отметим, что в~рамках приоритетов  

на\-уч\-но-тех\-но\-ло\-ги\-че\-ско\-го развития России активно работают 

научные коллективы Федерального исследовательского центра 

<<Информатика и~управление>> РАН (ФИЦ ИУ РАН), который 

сосредоточил свои научные ресурсы и~компетенции на исследовании 

научных проблем, определяемых прежде всего первым приоритетом, на 

основе принципиального подхода, суть которого~--- синергетика теории 

и~практики, предполагающая практическую направленность каждого 

результата фундаментальных исследований~[4]. 

    

    Поскольку ФИЦ ИУ РАН выступил инициатором КНТП~---  

<<Искусственный интеллект как драйвер цифровой трансформации 

экономики России>>, в~\mbox{статье} дана предварительная оценка потенциальных 

рисков и~угроз применительно к~этой программе. 

    

\section{Методические подходы к~классификации угроз}



    Используя опыт анализа угроз и~уязвимостей сложных систем~[5, 6], 

можно отметить, что негативные воздействия на ход формирования 

и~реализации КНТП могут оказывать как внешние, так и~внутренние 

факторы на всех стадиях жизненного цикла (ЖЦ). 

    

    К \textbf{внутренним факторам} отнесем те, которые обусловливаются 

процессами организации выполнения КНТП. Так, основным внутренним 

фактором является деятельность системы субъектов процесса реализации 

КНТП, связанная с~ее разработкой и~реализацией. При этом воздействие 

этого фактора может быть как положительным, так и~отрицательным. 

    

    \textbf{Внешние факторы} находятся вне Программы и~могут 

оказывать также аналогичное воздействие. Так, в~<<процессе определения 

потребностей и~требований заинтересованной стороны>> нужно учитывать 

интересы сторон, которые выступают \textit{против системы} (внешняя 

угроза) или \textit{друг против друга} (внут\-рен\-няя угроза)~[7]. 

Если интересы заинтересованных сторон направлены друг против друга, но 

не выступают против системы, данный процесс предназначен для 

достижения согласия среди заинтересованных сторон с~целью установления 

общего множества приемлемых требований. 

    

    Противодействие тех, кто находится в~оппозиции к~сис\-те\-ме, парируется 

с~использованием процессов управ\-ле\-ния рисками, анализа угроз, процессами 

сис\-тем\-но\-го анализа или системных требований к~безопас\-ности, 

адаптируемости или стойкости. В~этом случае, когда сталкиваются 

с~противодействием, потребности заинтересованной стороны не 

удовлетворяются, а,~скорее, направляются таким способом, чтобы помочь 

обеспечить системные гарантии и~целостность.

    

    Таким образом, в~качестве \textbf{первого критерия} классификации 

угроз можно выделить негативное воздействие внешних и~внутренних 

факторов.

    

    \textbf{Второй критерий} классификации угроз определяет разделение 

их по стадиям и~процессам ЖЦ. В~соответствии с~[7], 

типичные стадии ЖЦ системы включают замысел, разработку, производство, 

применение, поддержку и~выведение из эксплуатации, которые реализуются 

регламентированными процессами ЖЦ. Применительно к~ранним стадиям 

ЖЦ, к~которым, как правило, относят замысел, разработку и~производство, 

подлежат анализу с~точки зрения оценки угроз следующие процессы:

    \begin{itemize}

\item определение системных требований;

\item определение потребностей и~требований заинтересованной стороны;

\item управление решениями;

\item приобретение;

\item верификация, валидация и~приемка. 

\end{itemize}



    \textbf{Третий критерий} служит для дифференциации угроз на основе 

анализа рисков их воздействия на Программу. Риски количественно 

оценивают угрозу с~точки зрения вероятности ее реализации и~возможного 

ущерба. Их анализ позволяет выделить из большого множества 

существующих угроз наиболее существенные. 

    

    На основе стандартного подхода можно определить четыре степени 

опасности угроз (рис.~1): 

    \begin{enumerate}[(1)]

\item угрозы с~приемлемым риском, зачастую устранимые организационными 

мерами;

\item угрозы с~минимизируемым риском~--- путем изменения Программы 

с~\mbox{целью} уменьшения вероятности риска;

\item угрозы с~передаваемым риском~--- путем страхования, что требует 

выделения дополнительных ресурсов;

\item неприемлемые угрозы, парируемые изменениями в~Программе, 

поз\-во\-ля\-ющи\-ми избежать этих угроз.

\end{enumerate}



    \begin{figure} %fig1

    \vspace*{1pt}

 \begin{center}

 \mbox{%

 \epsfxsize=125.782mm 

\epsfbox{zac-1.eps}

 }

 \end{center}

\vspace*{-11pt}

    \Caption{Дифференциация рисков Программы}

     \end{figure}

     

    

\section{Классификатор угроз}



    Классификатор сформирован на основе приведенных выше критериев и~положений~[7] и~приведен на рис.~2.

    

    Основные угрозы на стадии замысла и~формирования Программы могут 

быть обусловлены как внутренними, так и~внешними причинами. 

Внутренние причины~--- конфликт интересов сторон, некомпетентность 

заказчика и~исполнителя, отсутствие независимой экспертной оценки~--- 

могут привести к~некорректному целеполаганию. Целеполагание здесь 

понимается в~широком смысле и~включает не только формулировку целей 

и~задач Программы, но и~системные требования с~учетом интересов сторон, 

перечень целевых показателей, критерии, этапность и~элементы 

планирования. При этом формулируемые цели должны быть конкретными, 

измеримыми, достижимыми, ресурсообеспеченными и~привязанными 

к~времени.



\begin{figure} %fig2

 \vspace*{1pt}

 \begin{center}

 \mbox{%

 \epsfxsize=128mm 

\epsfbox{zac-2.eps}

 }

 \end{center}

\vspace*{-6pt}

\Caption{Классификатор основных групп угроз}

\end{figure}

    

    Учет внешних источников угроз позволяет на стадии замысла 

Программы ослабить риски невозможности выполнения требований 

и~реализации решений в~силу внешних ограничений (санкций) 

и~организационных проблем, связанных с~нор\-ма\-тив\-но-пра\-во\-вой базой 

и~отсутствием независимого контроля.

    

    Степень опасности угроз на данной стадии ЖЦ высокая, так как 

угрожает существованию Программы и~может привести к~функционально 

критичным для Программы недостаткам.

    

    Основные уязвимости на \textbf{стадии реализации} Программы во 

многом совпадают со стадией замысла за следующими исключениями. 

В~процессе управления техническими решениями может выясниться 

невозможность реализации важнейших функциональных требований в~силу 

невозможности приобретения зарубежных технологий, например из-за 

санкций, и~отсутствия отечественных аналогов (реализуется уязвимость 

\textit{технологическое отставание}) или в~силу существования 

\textit{нерешенных научно-технических проб\-лем}, которое может быть 

обусловлено или \textit{наличием общенаучной проб\-ле\-мы}, или  

\textit{на\-уч\-но-тех\-ни\-че\-ским отставанием}. Кроме того, на данной 

стадии ЖЦ Программы понижается степень опасности уязвимостей до 

средней и~умеренной, так как большинство из них устраняется выделением 

дополнительных ресурсов или организационными мерами. Важнейшая 

уязвимость на всех стадиях ЖЦ~--- \textit{отсутствие независимого 

экспертного обеспечения и~научного сопровождения}, что не позволяет 

осуществлять объективную оценку как требований к~разрабатываемым 

системам, так и~результатов работы исполнителей. Одной из особенностей 

угроз, связанных с~проблематикой научных исследований, является 

возможность получения отрицательного научного результата, что должно 

считаться несомненно серьезным, но принимаемым риском.

    

\vspace*{-3pt}

    

\section{Комплексная научно-техническая программа <<Искусственный интеллект как драйвер 

цифровой трансформации экономики России>> и~оценка 

потенциальных угроз ее невыполнения}



\vspace*{-3pt}



    Первым приоритетом в~Стратегии научно-технологического развития 

Российской Федерации определен <<переход к~передовым цифровым, 

интеллектуальным, производственным технологиям, роботизированным 

системам, новым материалам и~способам конструирования, создание систем 

обработки больших объемов данных, машинного обучения и~искусственного 

интеллекта>>~[1]. И~выделение такого приоритета вполне оправдано 

стратегическим курсом руководства страны на цифровую трансформацию 

общества. При этом особую значимость в~рамках первого приоритета 

приобретает проблематика искусственного интеллекта (ИИ), который сегодня по 

праву относится к~одному из мегатрендов развития информационных 

технологий цифровой экономики. Поэтому технологии ИИ становятся одним из доминирующих факторов развития 

современного общества. Об этом свидетельствует ряд фактов. 

    

     Так, по прогнозам известной компании Accenture (США)~[8], 

ИИ может привести к~созданию <<виртуальной 

рабочей силы>>; наибольший рост за счет ИИ

ожидается в~американской экономике (4,6\%), в~Финляндии (4,1\%) и~Великобритании (3,9\%)~[8, 9]. Президент США своим Указом 

<<О~сохранении американского лидерства в~области искусственного 

интеллекта>> определил комплекс мероприятий по интенсификации работ 

в~об\-ласти ИИ~[10]. Стратегии по развитию ИИ 

приняты в~десятках стран мира.

    

     В~России Указом Президента РФ от 10.10.2019 №\,490 также 

утверждена Национальная стратегия развития искусственного интеллекта на 

период до 2030~г.~\cite{9-zat}. 

    

     К сожалению, Россия не входит в~число ведущих стран мира в~области 

ИИ. Достаточно упомянуть о публикационной активности российских 

ученых за последние десять лет (2009--2019~гг.): в~среднем каждая 

четвертая публикация~--- американская, каждая пятая~--- китайская, 

а~российская~--- одна из ста~[11]. 

    

     В связи с~этим инициатива ФИЦ ИУ РАН по разработке КНТП 

<<Искусственный интеллект как драйвер цифровой трансформации 

экономики России>> оказалась безусловно актуальной. Так, научными 

коллективами Центра в~рамках Совета по приоритетным направлениям 

Стратегии  

на\-уч\-но-тех\-но\-ло\-ги\-че\-ско\-го развития России в~октябре 2018~г.\ был 

разработан и~представлен проект Концепции комплексной  

на\-уч\-но-тех\-ни\-че\-ской программы <<Искусственный интеллект как 

драйвер цифровой трансформации экономики России>>, а~в~мае 2019~г.~--- 

проект комплексной программы. Оба документа были одобрены Советом 

(председатель~--- академик РАН И.\,А.~Каляев). 

    

     В КНТП предусмотрен комплекс мероприятий, охватывающих все 

стадии ЖЦ (фундаментальные исследования, разработка 

базовых технологий ИИ, создание инструментов 

и~ап\-па\-рат\-но-про\-грам\-мных средств ИИ, 

внед\-ре\-ние технологий ИИ в~различные сферы 

циф\-ро\-вой экономики, подготовка кадров в~об\-ласти ИИ). Планируется создание и~развитие 13~базовых технологий ИИ, 

а~также разработка и~внедрение более 60~прикладных технологий ИИ, 

8~университетских программ подготовки специалистов, создание не менее 

5~на\-уч\-но-об\-ра\-зо\-ва\-тель\-ных центров мирового уровня в~области 

ИИ, модели правового регулирования интеллектуальных технологий для 

баланса между сохранением фундаментальных общественных и~правовых 

ценностей, а~также ап\-па\-рат\-но-про\-грам\-мной инфраструктуры 

и~платформы поддержки решения задач ИИ, в~том 

числе специализированных суперкомпьютерных центров для решения задач 

в~об\-ласти ИИ %искусственного интеллекта 

и~глубокого машинного обучения.

    

     Важно, что Программа предусматривает комплекс работ, 

охватывающих все стадии ЖЦ сис\-тем, включая 

фундаментальные и~прикладные исследования, разработку и~внедрение 

технологий и~конкретных комплексов, а~также подготовку кадров. При этом 

сама Программа является, по существу, системой, развивающейся во времени 

и~включающей в~общем случае стадии замысла, формирования 

и~согласования требований, а также их последующей реализации. Таким 

образом, КНТП, одобренная Советом по приоритетным направлениям 

Стратегии на\-уч\-но-тех\-но\-ло\-ги\-че\-ско\-го развития России, стала, по 

существу, основой для постановки на\-уч\-но-прак\-ти\-че\-ских задач 

в~рамках реализации утвержденной Президентом Стратегии. 

    

     Вместе с~тем уже сейчас можно говорить о сложности реализации 

КНТП, обусловленной прежде всего  

ор\-га\-ни\-за\-ци\-он\-но-ме\-то\-ди\-че\-ски\-ми проб\-ле\-ма\-ми, которые 

создают реальные риски либо невыполнения Программы, либо снижения 

эффективности запланированных в~ней мероприятий. 



\section{Оценка потенциальных угроз невыполнения комплексной научно-технической программы 

<<Искусственный интеллект как~драйвер цифровой 

трансформации экономики России>>}



    В этом разделе сделана попытка формализованной оценки рисков 

невыполнения КНТП с~учетом предложенной выше классификации 

и~дифференциации угроз по четырем степеням. 

    

    \textbf{По первому критерию}, который включает внутренние 

и~внешние факторы. Внутренние факторы определяются содержанием 

КНТП, а также системой субъектов для ее выполнения. 

    

    \textbf{Вероятность ошибки целеполагания~--- низкая}. Такая оценка 

обусловлена тем, что целью КНТП ставится создание технологий 

ИИ, обеспечивающих повышение 

производительности труда, создание новых видов продукции, рост 

экономического благосостояния, продолжительности и~качества жизни 

российского общества, укрепление национальной безопасности и~достижение 

национальных приоритетов на международной арене. Другими словами, 

целеполагание КНТП в~полной мере соответствует первому приоритету 

Стратегии на\-уч\-но-тех\-но\-ло\-ги\-че\-ско\-го развития и~принятым 

документам по цифровой трансформации общества. 

    

    \textbf{Вероятность ошибки реализации~--- умеренная}. 

В~настоящее время затруднительно сделать объективную оценку, так как 

еще не сформирована система субъектов, определяющая организационные, 

методические и~финансовые процессы выполнения Программы. При этом 

ключевым субъектом в~этой системе становится ответственный исполнитель, 

на которого возлагается организация и~координация работ, оперативное 

управление и~контроль в~рамках Программы. Однако ответственный 

исполнитель до сих пор не определен. 

    

    В этих условиях с~определенной степенью условности примем 

вероятность ошибки реализации как умеренную. 

    

    Гораздо сложнее обстоят дела с~оценкой \textbf{внешних факторов}. 

К~ним следует отнести определенный конфликт между КНТП 

и~утвержденной Национальной стратегией развития искусственного 

интеллекта~\cite{9-zat}. Кратко о сути этого конфликта~\cite{13-zat}. 

    

    Стратегия разработана ПАО <<Сбербанк>> без привлечения научных 

организаций, и~среди определенных в~ней принципов, к~сожалению, 

отсутствует акцент на интенсивное проведение научных исследований 

в~сфере ИИ. Так, в~качестве основной цели развития 

ИИ в~России определено <<\textit{формирование 

самостоятельной индустрии искусственного интеллекта, массовое 

использование технологий искусственного интеллекта в~национальной 

экономике и~обеспечение международного лидерства по отдельным 

направлениям применения технологий искусственного 

интеллекта}>>~\cite{9-zat}. Более того, при подготовке этого документа не 

были в~полной мере учтены замечания ведущих российских ученых 

в~области ИИ, не нашел отражения тот факт, что в~России уже более~30~лет 

существует профессиональное сообщество и~научные школы в~области ИИ, 

вполне достойно представленные как на европейском, так и~на мировом 

уровне. 

    

    Тем не менее этот конфликт не является принципиальным, поскольку 

оба документа направлены на достижение целей, связанных со скорейшим 

продвижением методов и~технологий ИИ в~России. 

    

    Вместе с~тем среди внешних факторов, которые \textbf{позитивно} 

влияют на выполнение КНТП, необходимо назвать:

    \begin{itemize}

\item государственный курс на цифровую трансформацию общества, 

закрепленный в~ряде нормативных документов, в~том числе в~Национальных 

проектах <<Цифровая экономика>>, <<Наука>>, <<Инфраструктура>>;

\item достаточно высокий уровень базового физико-математического 

образования в~России (МГУ им.\ М.\,В.~Ломоносова, Физтех, МИФИ и~др.), 

научные школы в~области математики и~естественных наук (ФИЦ ИУ РАН, 

ИПМ РАН им.\ М.\,В.~Келдыша, ИПУ РАН им.\ Трапезникова и~др.); 

\item наличие приемлемой базовой ИКТ-инфраструктуры (распространение 

4G-се\-тей, большое число функциональных смартфонов, доступность сети 

Интернет);

\item потенциально огромный внутренний рынок для широкого 

использования технологий ИИ в~промышленности, добывающих отраслях, 

сельском хозяйстве, энергетике, транспорте. 

\end{itemize}

    

    Однако существует и~целый ряд внешних факторов, которые 

\textbf{негативно} будут влиять на выполнение КНТП. Выделим 

следующие: 

    \begin{itemize}

\item практическое отсутствие отечественных производителей аппаратного 

обеспечения для технологий ИИ;

\item отсутствие отечественного базового программного обеспечения для 

реализации методов и~алгоритмов ИИ;

\item слабый уровень развития облачных систем и~технологий;

\item недостаточный уровень специализированных вычислительных 

мощностей;

\item недостаточное число высококвалифицированных кадров в~сфере ИИ 

(ученых, инженеров, разработчиков, в~том числе в~связи с~отъездом 

талантливых специалистов из страны);

\item несовершенство нормативно-правовой базы в~области применения ИИ в~различных сферах, включая защиту прав граждан;

\item недостаточный уровень общественного доверия и~общей готовности 

населения к~широкому применению технологий ИИ, особенно в~сферах 

опасного производства, транспорта, здравоохранения и~других областях с~высокими рисками;

\item практическое отсутствие мер, направленных на создание 

благоприятных условий для внедрения российских средств ИИ на мировом 

уровне в~условиях возрастающей конкурентности в~этой области;\\[-14pt]

\item огромные потенциальные возможности технологий ИИ для их 

применения в~противоправных целях (кибермошенничество, включая кражи, 

терроризм и~т.\,п.);\\[-14pt] 

\item санкции (как существующие, так и~возможные новые) на запрет доступа 

к~современным технологиям и~средствам, зарубежным рынкам, на поставку 

аппаратного и~программного обеспечения, специфических средств 

жизнеобеспечения и~др.).\\[-13pt] 

    \end{itemize}

    

    Приведенные факторы создают целый ряд барьеров для разработки 

и~внед\-ре\-ния решений ИИ в~России и~реализации КНТП.

    

    Наряду с~приведенными выше внешними факторами, имеющими 

негативный характер, следует отдельно отметить такой принципиально 

важный вопрос, как финансирование КНТП. Так, финансовое обеспечение 

реализации комплексной программы планируется осуществлять 

ответственными исполнителями, соисполнителями и~участниками КНТП. 

При этом в~качестве источников финансирования определены: федеральный 

бюджет, бюджеты субъектов Российской Федерации, местные бюджеты 

и~средства внебюджетных источников. Другими словами, целевое 

финансирование КНТП не предусматривается. Этот фактор создает 

наибольшие риски невыполнения КНТП. 

    

    Таким образом, внешние факторы по совокупности позволяют оценить 

\textbf{вероятность невыполнения КНТП как высокую}. 

    

    Каковы пути нейтрализации выявленных угроз и~рисков? Очевидно, что 

их можно представить в~виде двух групп. 

    

    \textbf{Первая группа} включает комплекс мероприятий, который 

должен кардинально изменить отношение в~целом к~науке в~России. Без 

этого обеспечить ведущее положение страны в~области ИИ невозможно. По 

данным Института статистических исследований и~экономики знаний НИУ 

ВШЭ~\cite{12-zat}, Россия находится на 10-м месте в~рейтинге ведущих 

стран мира по величине затрат на научные исследования и~разработки. 

Однако по доле затрат на науку в~валовом внутреннем продукте (ВВП) наша 

страна находится на 34-м месте (с~показателем~1,1\%), по величине затрат 

в~расчете на одного исследователя~--- на 47-м месте (93~тыс.\ долл.\ США в~год), а по 

численности исследователей в~расчете на 10~тыс. занятых в~экономике~--- на 

34-м месте. Для сравнения: лидеры мировой экономики США и~Китай 

выделяют на финансирование науки из своих огромных ВВП 

соответственно~2,74\% (11-е место) и~2,12\% (15-е место).  

А~стра\-на\-ми-ли\-де\-ра\-ми, руководство которых осознало роль 

и~значимость научных исследований в~развитии экономики и~общества 

и~поэтому обеспечивает опережающее их финансирование, стали Израиль 

(4,45\%), Финляндия (4,30\%), Южная Корея (4,24\%), Швеция (4,12\%), 

Япония (3,92\%). 

%   

 Отсюда следует, что необходимо обеспечить повышение уровня 

финансирования научных исследований \textbf{не менее чем в~2~раза (до 

2,2\%--2,5\% от ВВП)}. 



\pagebreak

    

    Необходимо в~корне изменить систему оценки научных исследований. 

Сегодняшняя оценка по числу публикаций в~изданиях, цитируемых 

в~зарубежных базах (WoS, Scopus и~др.), не позволяет реально оценить 

состояние российской науки и~не способствует повышению 

результативности научных исследований. Более того, в~этом году эта оценка 

стала еще более абстрактно формализованной~--- в~числе баллов, которые 

являются весьма сложной функцией от реальных научных результатов. 

В~таких подходах кроется принципиальное смешение понятий: оценка 

научной деятельности конкретных исследователей и~оценка 

результативности научной деятельности научных коллективов 

и~организаций. Поэтому необходимо вернуться к~системе оценки научных 

организаций по достигнутым научным результатам, реализуемым 

в~конкретных областях промышленности и~экономики. 

    

    \textbf{Вторая группа} включает проведение конкретных 

организационных мероприятий применительно к~КНТП: скорейшее 

определение всех субъектов реализации КНТП, прежде всего ответственного 

исполнителя, выработку единой согласованной позиции в~рамках системы 

субъектов реализации КНТП, а также утвержденной Национальной 

стратегии. Совершенно очевидно, что мероприятия в~рамках этих двух 

документов должны дополнять друг друга с~акцентами: 

    \begin{itemize}

\item в~КНТП~--- на научную обоснованность на основе фундаментальных 

исследований и~создание базовых отечественных наукоемких технологий 

ИИ, а~также целевую подготовку 

высококвалифицированных специалистов, способных эффективно применять 

технологии ИИ; 

\item в~Национальной стратегии~--- на создание прикладных технологий на 

основе базовых, массовое производство изделий, реализующих эти 

технологии, а~также вопросы их внедрения и~применения в~различных 

областях. 

\end{itemize}



    Очевидно, что и~вопросы определения порядка и~источников 

финансирования должны быть дифференцированы по разделам программы:

    \begin{itemize}

\item фундаментальные исследования и~подготовка кадров~--- Минобрнауки 

России; 

\item базовые технологии~--- Минобрнауки совместно с~заинтересованными 

ведомствами; 

\item прикладные технологии и~их внедрение~--- конкретными 

заинтересованными заказчиками, включая государственные и~коммерческие 

организации. 

\end{itemize}

 

\section{Заключение }



    Комплексные научно-технические программы исследований и разработок

являются актуальным инструментом реализации Стратегии 

научно-тех\-но\-ло\-ги\-че\-ско\-го развития России в~рамках принятых приоритетов. 

Пред\-ла\-га\-емые в~\mbox{статье} методические подходы к~классификации рисков 

и~угроз на основе их систематизации позволяют дать качественную оценку 

возможности выполнения конкретных КНТП, а~также получить некоторое 

цельное представление об эффективности запланированного в~КНТП 

комплекса мероприятий. 

    

    Проведенная оценка потенциальных рисков и~угроз применительно 

к~КНТП <<Искусственный интеллект как драйвер цифровой трансформации 

экономики России>>, разработанной ФИЦ ИУ РАН, показала, что 

с~достаточно высокой вероятностью эта Программа может быть не 

выполнена. 

    %

    При этом основные риски обусловлены внешними факторами: 

недостаточным стимулированием научных исследований и~нерешенностью 

организационных вопросов.



 Сегодня актуальной становится синергетика 

теории и~практики~--- результатов фундаментальных исследований и~их 

реализации в~конкретных системах и~технологиях. Для этого необходимо 

объединить и~науку (фундаментальную, прикладную, военную), 

и~технологии, и~промышленное производство. 

    

    Федеральный исследовательский центр <<Информатика и~управ\-ле\-ние>> Российской академии наук

     обладает большим научным потенциалом в~области 

компьютерных наук, а также уникальным синтезом результатов глубоких 

фундаментальных исследований в~области информатики и~огромного  

на\-уч\-но-прак\-ти\-че\-ско\-го опыта разработки, внедрения, модернизации 

и~сопровождения ин\-фор\-ма\-ци\-он\-но-те\-ле\-ком\-му\-ни\-ка\-ци\-он\-ных сис\-тем 

в~интересах органов государственной власти на различных уровнях. 



{\small\frenchspacing

{%\baselineskip=10.8pt

\addcontentsline{toc}{section}{Литература}

\begin{thebibliography}{99}



\bibitem{1-zat}

О Стратегии научно-технологического развития Российской Федерации: Указ Президента 

Российской Федерации от 01.12.2016 №\,642.

{\sf 

http://static.\linebreak kremlin.ru/media/events/files/ru/uZiATIOJiq5tZsJgqcZLY9YyL8PWTXQb. pdf}.

\bibitem{2-zat}

Об утверждении Правил разработки, утверждения, реализации, корректировки 

и~завершения комплексных на\-уч\-но-тех\-ни\-че\-ских программ полного инновационного цикла 

и~комплексных на\-уч\-но-тех\-ни\-че\-ских проектов полного инновационного цикла в~целях 

обеспечения реализации приоритетов научно-технологического развития Российской 

Федерации: Постановление Правительства РФ от 19.02.2019 №\,162.

{\sf 

http://static.government.ru/media/files/TX7NZe8Am8Ovkf0UEgDVgliHIktbAUK2.\linebreak pdf}.

\bibitem{3-zat}

Рейтинг экономики развитых стран мира~/\!\!/ VisaSam.ru, 21.06.2020. {\sf 

https://\linebreak visasam.ru/emigration/vybor/ekonomika-stran-mira-2.html}.

\bibitem{4-zat}

\Au{Зацаринный А.\,А.} Научные исследования в~интересах цифровой трансформации 

общества в~условиях первого приоритета научно-технологического развития России~/\!\!/ 

Проектирование будущего. Проблемы цифровой реальности~/ Под ред. 

Г.\,Г.~Малинецкого.~--- М.: ИПМ им.\ М.\,В.~Келдыша, 2020 (в печати).

\bibitem{5-zat}

\Au{Сучков А.\,П.} Классификация уязвимостей интегрированных систем управления на 

ранних стадиях жизненного цикла~/\!\!/ Системы и~средства информатики, 2017. Т.~27. 

№\,4. С.~132--143.

\bibitem{6-zat}

\Au{Сучков А.\,П.} Процессная модель модернизации и~развития информационных систем 

на всех стадиях жизненного цикла~/\!\!/ Системы и~средства информатики, 2018. Т.~28. 

№\,4. С.~84--95 .

\bibitem{7-zat}

%ГОСТ ИСО/МЭК/ИИРЭ 24748-1. %письмо от ОБФ 13.10.20

ГОСТ Р 56923-2016/ISO/IEC TR 24748-3:2011. 

Информационные технологии. Системная и~программная инженерия. 

Управление жизненным циклом. Часть~3. Руководство по применению 

ИСО/МЭК 12207 (Процессы жизненного цикла программных средств).~--- 

М.: Стандартинформ, 2016. 217~с.



\bibitem{8-zat}

Accenture: Искусственный интеллект ускорит ежегодные темпы экономического рос\-та 

к~2035~году~/\!\!/ Inc./Новости, 01.02.2017. {\sf 

 https://incrussia.ru/news/accenture-iskusstvennyy-intellekt-uskorit-ezhegodnye-tempy-ekonomicheskogo-rosta-k-2035-godu}.

\bibitem{9-zat}

О развитии искусственного интеллекта в~Российской Федерации: Указ Президента 

Российской Федерации от 10.10.2019 №\,490.

\bibitem{10-zat}

Executive Order on Maintaining American Leadership in Artificial Intelligence. 11.02.2019. {\sf 

https://www.whitehouse.gov/presidential-actions/executive-order-maintaining-american-leadership-artificial-intelligence/}. 



\bibitem{11-zat}

\Au{Городникова Н.\,В., Гохберг~Л.\,М., Дитковский~К.\,А. и~др.} Наука. Технологии. 

Инновации: 2019: Краткий статистический сборник.~--- М.: НИУ ВШЭ, 2019. 84~с.



\bibitem{13-zat}

\Au{Устинова А.} Ученые призвали доработать нацстратегию по искусственному 

интеллекту~/\!\!/ ComNews, 28.06.2019. {\sf  

https://www.comnews.ru/content/120485/2019-06-28/uchenye-prizvali-dorabotat-nacstrategiyu-po-iskusstvennomu-intellektu.}



\bibitem{12-zat}

Расходы на науку: топ-10 стран мира~/\!\!/ Рамблер/Новости, 25.07.2018. {\sf 

https:// news.rambler.ru/other/40408588-rashody-na-nauku-top-10-stran-mira/}.

    \end{thebibliography}

} }





\vspace*{-3pt}



\hfill{\small\textit{Поступила в~редакцию 02.04.20}}









\vspace*{6pt}



\hrule



\vspace*{2pt}



%\newpage



%\vspace*{-28pt}



\hrule



\vspace*{8pt}



\def\tit{THREATS AND RISKS OF~IMPLEMENTING COMPLEX SCIENTIFIC AND~TECHNICAL 

PROGRAMS\\ WITHIN THE~PRIORITIES OF~THE~RUSSIAN SCIENCE AND~TECHNOLOGY 

DEVELOPMENT STRATEGY}



\def\titkol{Threats and risks of~implementing CSTP} %complex scientific and~technical programs 

%within the~priorities of~the~Russian science and~technology development strategy}



\def\aut{A.\,A.~Zatsarinny and~A.\,P.~Suchkov}



\def\autkol{A.\,A.~Zatsarinny and~A.\,P.~Suchkov}



\titel{\tit}{\aut}{\autkol}{\titkol}



\vspace*{-11pt}





\noindent

Institute of Informatics Problems, Federal Research Center 

``Computer Science and Control'' of the Russian 

Academy of Sciences, 44-2~Vavilov Str., Moscow 119133, Russian Federation



\def\leftfootline{\small{\textbf{\thepage}

\hfill Sistemy i Sredstva Informatiki~---

Systems and Means of Informatics 2020 vol~30 no\,3}

}%

 \def\rightfootline{\small{Sistemy i Sredstva Informatiki~---

 Systems and Means of Informatics 2020 vol~30 no\,3

\hfill \textbf{\thepage}}}



\vspace*{3pt}



\Abste{The article deals with the actual problems of implementing the strategy of scientific and 

technological development of Russia in the framework of accepted priorities based on complex scientific 

and technical research and development programs (CSTP), primarily within the first priority. 

Methodological approaches to classification of risks and threats based on their systematization, which 

hinder implementation of the adopted CSTP, are presented. A threat classifier is

proposed that allows one to get a complete picture of effectiveness of the planned

set of measures in the CSTP. An assessment of 

potential risks and threats in relation to the ``Artificial intelligence as a driver of digital transformation of 

the Russian economy,'' developed by the Federal Research Center ``Computer\linebreak\vspace*{-12pt}}



\Abstend{Science and Control'' of the 

Russian Academy of Sciences, is given. It is shown that the main risks are caused by insufficient 

stimulation of scientific research and low demand for innovative scientific results.}



\KWE{complex scientific and technical program; strategy of scientific and technological development; 

classification of risks and threats; artificial intelligence; scientific research}







\DOI{10.14357\!/\!08696527200309}



%\vspace*{-24pt}



 \Ack

\noindent

The paper was partially supported by the Russian Foundation for Basic Research (project 18-29-03091).



 



\vspace*{-12pt}



\renewcommand{\bibname}{\protect\rmfamily References}

%\renewcommand{\bibname}{\large\protect\rm References}



{\small\frenchspacing

{%\baselineskip=10.8pt

\addcontentsline{toc}{section}{References}

\begin{thebibliography}{99}

\bibitem{1-zat-1}

O strategii nauchno-tekhnologicheskogo razvitiya Rossiyskoy Federatsii: ukaz Prezidenta ot 01.12.2016 

No.\,642 [About strategy of scientific and technological development of the Russian Federation. 

Presidential Decree No.\,642 dated 01.12.2016]. Available at: {\sf 

http://static.kremlin.ru/media/events/files/ru/uZiATIOJiq5tZsJgqcZLY9YyL8PW\linebreak TXQb.pdf} 

(accessed August~14, 2020).

\bibitem{2-zat-1}

Ob utverzhdenii Pravil razrabotki, utverzhdeniya, realizatsii, korrektirovki i~zaversheniya kompleksnykh 

nauchno-tekhnicheskikh programm polnogo innovatsionnogo tsikla i~kompleksnykh  

nauchno-tekhnicheskikh proektov polnogo innovatsionnogo tsikla v~tselyakh obespecheniya realizatsii 

prioritetov nauchno-tekhnologicheskogo razvitiya Rossiyskoy Federatsii: Postanovlenie Pravitel'stva RF 

ot 19.02.2019 No.\,162 [On the approval of the Rules for the development, approval, 

implementation, adjustment, and completion of complex scientific and technical programs of a full 

innovation cycle and complex scientific and technical projects of a full innovation cycle in order to ensure 

the implementation of the priorities of scientific and technological development of the Russian 

Federation: Government decree dated February~16, 2019 No.\,162]. Available at: {\sf 

http://static.government.ru/media/files/TX7NZe8Am8Ovkf0UEgDVgliHIktbAUK2.\linebreak pdf} (accessed 

August~14, 2020).

\bibitem{3-zat-1}

Reyting ekonomiki razvitykh stran mira [World developed countries economy rating]. 

21.06.2020. VisaSam.ru.

Available at: {\sf 

https://visasam.ru/emigration/ vybor/ekonomika-stran-mira-2.html} (accessed August~14, 2020).

\bibitem{4-zat-1}

\Aue{Zatsarinny, A.\,A.} 2020 (in press). Nauchnye issledovaniya v~interesakh tsifrovoy transformatsii 

obshchestva v~usloviyakh pervogo prioriteta nauchno-tekhnologicheskogo razvitiya Rossii [Scientific 

research in the interests of the digital transformation of society in the context of the first priority of the 

scientific and technological development of Russia]. \textit{Proektirovanie budushchego. Problemy 

tsifrovoy real'nosti} [The design of the future. Problems of digital reality]. Ed.~G.\,G.~Malinetskiy. Moscow: IPM. 

\bibitem{5-zat-1}

\Aue{Suchkov, A.\,P.} 2017. Klassifikatsiya uyazvimostey integrirovannykh sistem upravleniya na 

rannikh stadiyakh zhiznennogo tsikla [Classification of vulnerabilities of integrated management systems 

at the early stages of the lifecycle]. \textit{Sistemy i~Sredstva Informatiki~--- Systems and Means of 

Informatics} 27(4):132--143.

\bibitem{6-zat-1}

\Aue{Suchkov, A.\,P.} 2018. Protsessnaya model' modernizatsii i~razvitiya informatsionnykh sistem na 

vsekh stadiyakh zhiznennogo tsikla [Process model of modernization and development of information 

systems at all stages of the life cycle]. \textit{Sistemy i~Sredstva Informatiki~--- Systems and Means of 

Informatics} 28(4):84--95. 

\bibitem{7-zat-1}

GOST R 56923-2016/ISO/IEC TR 24748-3:2011. 2016. Informatsionnye tekhnologii. 

Sistemnaya i programmnaya inzheneriya. Upravlenie zhiznennym tsiklom. 

Chast'~3. Rukovodstvo po primeneniyu ISO/MEK 12207 (Protsessy zhiznennogo tsikla programmnykh 

sredstv) [Information technologies. Systems and software engineering. Life cycle management. Part~3. Guide to the 

application of ISO/IEC 12207 (Software life cycle processes)]. Moscow: Standardinform Publs. 217~p. 



\bibitem{8-zat-1}

Accenture: Iskusstvennyy intellekt uskorit ezhegodnye tempy ekonomicheskogo rosta k~2035 godu [Accenture: Artificial 

intelligence will accelerate the annual rate of economic growth by 2035]. 01.02.2017. Inc.\!/\!Novosti [Inc.\!/\!News]. Available at: {\sf 

https://incrussia.ru/news/accenture-iskusstvennyy-intellekt-uskorit-ezhegodnye-tempy-ekonomicheskogo-rosta-k-2035-godu/} 

(accessed August~14, 2020).

\bibitem{9-zat-1}

O razvitii iskusstvennogo intellekta v Rossiyskoy Federatsii: ukaz Prezidenta ot 10.10.2019 No.\,490 

[About strategy of scientific and technological development of the Russian Federation. Presidential 

Decree No.\,490 dated 10.10.2019]. Available at: {\sf 

http://static.kremlin.ru/media/events/files/ru/AH4x6HgKWANwVtMOfPDhcbRpv\linebreak d1HCCsv.pdf} 

(accessed August~14, 2020).

\bibitem{10-zat-1}

Executive Order on Maintaining American Leadership in Artificial Intelligence. 11.02.2019. Available at: {\sf 

https://www.whitehouse.gov/presidential-actions/executive-order-maintaining-american-leadership-artificial-intelligence/} 

(accessed August~14, 2020).

\bibitem{11-zat-1}

\Aue{Gorodnikova, N.\,V., L.\,M. Gokhberg, K.\,A.~Ditkovskiy, \textit{et. al.}} 2019. \textit{Nauka. 

Tehnologii. Innovatsii: Kratkiy statisticheskiy sbornik} [The science. Technologies. Innovations: A~brief 

statistical compendium]. Moscow: NIU HSE. 84~p.



\bibitem{13-zat-1}

\Aue{Ustinova, A.} 28.06.2019. Uchenye prizvali dorabotat' natsstrategiyu po iskusstvennomu intellektu 

[Scientists urged to finalize the national strategy for artificial intelligence]. \textit{ComNews}. Available 

at: {\sf  

https://www.comnews.ru/content/120485/2019-06-28/uchenye-prizvali-dorabotat-nacstrategiyu-po-iskusstvennomu-intellektu} 

(accessed August~14, 2020).



\bibitem{12-zat-1}

Raskhody na nauku: top-10 stran mira [Spending on science: Top~10 countries in the world].  23.07.2018.

Rambler\!/\!Novosti [Rambler\!/\!News]. Available 

at: {\sf https://news.rambler.ru/other/40408588-rashody-na-nauku-top-10-stran-mira/} (accessed 

August~14, 2020).

\end{thebibliography}

} }



\vspace*{-3pt}



\hfill{\small\textit{Received April 2, 2020}} 



\vspace*{-16pt}



\Contr



\vspace*{-4pt}



\noindent

\textbf{Zatsarinny Alexander A.} (b.\ 1951)~--- Doctor of Science in technology, professor, Deputy 

Director, Federal Research Center ``Computer Science and Control'' of the Russian Academy of Sciences 

(FRC CSC RAS); principal scientist, Institute of Informatics Problems, FRC CSC RAS; 44-2~Vavilov 

Str., Moscow 119333, Russian Federation; \mbox{AZatsarinny@ipiran.ru}



\vspace*{3pt}



\noindent

\textbf{Suchkov Alexander P.} (b.\ 1954)~--- Doctor of Science in technology, leading scientist, Institute 

of Informatics Problems, Federal Research Center ``Computer Science and Control'' of the Russian 

Academy of Sciences, 44-2~Vavilov Str., Moscow 119333, Russian Federation; 

\mbox{ASuchkov@frccsc.ru}





\label{end\stat}



\renewcommand{\bibname}{\protect\rm Литература} 

